\section{Cómo elegir entre STaR, GRPO y DAPO}

\videofigure{figures/method-selection.jpg}{Distintas escalas de datos, longitudes de trayectoria y presupuestos de cómputo corresponden a distintos métodos de entrenamiento.}{01:02:26--01:03:34}

\begin{table}[H]
\centering
\small
\begin{tabular}{>{\raggedright\arraybackslash}p{2.2cm}>{\raggedright\arraybackslash}p{3.6cm}>{\raggedright\arraybackslash}p{3.6cm}>{\raggedright\arraybackslash}p{4.0cm}}
\toprule
\textbf{Método} & \textbf{Condiciones de aplicación} & \textbf{Ventaja principal} & \textbf{Riesgo principal} \\
\midrule
STaR & Pocos rationale examples, respuesta verificable & Simple y eficiente, datos autogenerados & Sesgo de filtrado final, rationalization mismatch \\
GRPO & Se puede muestrear en línea un grupo de candidatos, la reward es calculable & Sin critic, sistema más ligero & Sin gradiente en grupos todo-correcto o todo-incorrecto, colapso de entropía \\
DAPO & CoT largo, RL en línea a gran escala & Escalamiento estable, mayor tasa de cómputo efectivo & Implementación compleja, hiperparámetros acoplados, requiere monitoreo intensivo \\
\bottomrule
\end{tabular}
\end{table}

\subsection{Por qué las tres métricas pueden diferir}

\videofigure{figures/metrics.jpg}{El RL puede cambiar simultáneamente pass@1, majority@$k$ y pass@$k$.}{01:03:34--01:06:20}

\begin{itemize}
    \item \textbf{pass@1}: un único muestreo se vuelve más correcto;
    \item \textbf{majority@$k$}: la respuesta correcta se convierte en la moda más común;
    \item \textbf{pass@$k$}: cobertura diversa del conjunto de candidatos.
\end{itemize}

El RL suele mejorar de manera notable majority@$k$, pero no necesariamente pass@$k$. Una explicación es que la policy concentra la probabilidad en los patrones correctos ya existentes, lo que aumenta la consistencia pero reduce la exploración diversa; esta es también la tensión entre «volverse más fuerte en el entrenamiento» y «ampliar la cobertura».

\subsection{Preguntas abiertas}

\videofigure{figures/open-problems.jpg}{El escalamiento del RL aún carece de una teoría completa sobre la consistencia mayoritaria, la cobertura, la diversidad y la conducta de autoevaluación.}{01:06:20--01:10:46}

\begin{importantbox}{La pregunta que más vale la pena seguir}
¿Por qué el RL hace más común la respuesta correcta sin necesariamente permitir que el modelo resuelva más problemas distintos? Responder esta pregunta requiere analizar a la vez la forma de la reward, la entropía, la distribución de dificultad de los datos y los puntos ciegos del verifier, en lugar de mirar solo el puntaje promedio del benchmark.
\end{importantbox}

\subsection{Resumen de la sección}

STaR es adecuado para el bootstrapping de bajo costo, GRPO es adecuado para la optimización relativa dentro del grupo sin critic, y DAPO es adecuado para el entrenamiento estable a gran escala con reasoning largo. La elección del método debe basarse en la verificabilidad de la retroalimentación, el costo del muestreo en línea y la longitud de la trayectoria.
